% =====================================================================
% Figure 3 — §VII-D IICF round trajectories (supplementary).
%
% Drop-in include for the AAAI 2027 paper supplementary appendix.
% Generated by safe-rl-2027/experiments/diagnostics/D6_iterative_critic/
% plot_d6.py from 5 seeds × {4, 6} rounds of round.json files.
% Native canvas is 8.8 × 3.4 in — two-column-figure width.
%
% Usage in §VII-D body:
%
%   % =====================================================================
% Figure 3 — §VII-D IICF round trajectories (supplementary).
%
% Drop-in include for the AAAI 2027 paper supplementary appendix.
% Generated by safe-rl-2027/experiments/diagnostics/D6_iterative_critic/
% plot_d6.py from 5 seeds × {4, 6} rounds of round.json files.
% Native canvas is 8.8 × 3.4 in — two-column-figure width.
%
% Usage in §VII-D body:
%
%   % =====================================================================
% Figure 3 — §VII-D IICF round trajectories (supplementary).
%
% Drop-in include for the AAAI 2027 paper supplementary appendix.
% Generated by safe-rl-2027/experiments/diagnostics/D6_iterative_critic/
% plot_d6.py from 5 seeds × {4, 6} rounds of round.json files.
% Native canvas is 8.8 × 3.4 in — two-column-figure width.
%
% Usage in §VII-D body:
%
%   % =====================================================================
% Figure 3 — §VII-D IICF round trajectories (supplementary).
%
% Drop-in include for the AAAI 2027 paper supplementary appendix.
% Generated by safe-rl-2027/experiments/diagnostics/D6_iterative_critic/
% plot_d6.py from 5 seeds × {4, 6} rounds of round.json files.
% Native canvas is 8.8 × 3.4 in — two-column-figure width.
%
% Usage in §VII-D body:
%
%   \input{figures/fig_vii_d_iicf_rounds/fig_vii_d_iicf_rounds.tex}
%
% Cross-references: §VII-D §D.3 (results table) and §D.4
% (diagnosis — two-fixed-point interpretation).
% =====================================================================
\begin{figure*}[t]
  \centering
  \includegraphics[width=\textwidth]%
    {figures/fig_vii_d_iicf_rounds/fig_vii_d_iicf_rounds.pdf}
  \caption{%
    \textbf{IICF round trajectories — the bimodal final outcome
    persists across $R \in \{4, 6\}$.}
    Per-seed episode cost over the outer-loop rounds for the v3b
    configuration (window depth $W\!=\!2$, per-round critic reset)
    at $R\!=\!4$ (left) and the same configuration extended to
    $R\!=\!6$ (right).  Each line is one of five seeds; per-round
    eval uses 20 episodes and is noisy.  The stars at the right
    edge mark the final \emph{deterministic 50-episode eval},
    coloured by the final verdict (green: $\bar C \le 5$, red:
    $\bar C > 5$).
    Three of five seeds in each run land cleanly inside the GO band
    while two settle at $\bar C\!\ge\!19$ — the same 3/5 split
    occurs in both panels, with different seed identities.
    Extending from $R\!=\!4$ to $R\!=\!6$ does not rescue the
    NO-GO seeds: more rounds do not change which basin a seed
    converges to.  This is the §D.4 ``two fixed points''
    diagnosis — the iteration converges, but to either of two
    attractors (safe versus over-conservative), determined by the
    initial PPO trajectory.%
  }
  \label{fig:vii_d_iicf_rounds}
\end{figure*}

%
% Cross-references: §VII-D §D.3 (results table) and §D.4
% (diagnosis — two-fixed-point interpretation).
% =====================================================================
\begin{figure*}[t]
  \centering
  \includegraphics[width=\textwidth]%
    {figures/fig_vii_d_iicf_rounds/fig_vii_d_iicf_rounds.pdf}
  \caption{%
    \textbf{IICF round trajectories — the bimodal final outcome
    persists across $R \in \{4, 6\}$.}
    Per-seed episode cost over the outer-loop rounds for the v3b
    configuration (window depth $W\!=\!2$, per-round critic reset)
    at $R\!=\!4$ (left) and the same configuration extended to
    $R\!=\!6$ (right).  Each line is one of five seeds; per-round
    eval uses 20 episodes and is noisy.  The stars at the right
    edge mark the final \emph{deterministic 50-episode eval},
    coloured by the final verdict (green: $\bar C \le 5$, red:
    $\bar C > 5$).
    Three of five seeds in each run land cleanly inside the GO band
    while two settle at $\bar C\!\ge\!19$ — the same 3/5 split
    occurs in both panels, with different seed identities.
    Extending from $R\!=\!4$ to $R\!=\!6$ does not rescue the
    NO-GO seeds: more rounds do not change which basin a seed
    converges to.  This is the §D.4 ``two fixed points''
    diagnosis — the iteration converges, but to either of two
    attractors (safe versus over-conservative), determined by the
    initial PPO trajectory.%
  }
  \label{fig:vii_d_iicf_rounds}
\end{figure*}

%
% Cross-references: §VII-D §D.3 (results table) and §D.4
% (diagnosis — two-fixed-point interpretation).
% =====================================================================
\begin{figure*}[t]
  \centering
  \includegraphics[width=\textwidth]%
    {figures/fig_vii_d_iicf_rounds/fig_vii_d_iicf_rounds.pdf}
  \caption{%
    \textbf{IICF round trajectories — the bimodal final outcome
    persists across $R \in \{4, 6\}$.}
    Per-seed episode cost over the outer-loop rounds for the v3b
    configuration (window depth $W\!=\!2$, per-round critic reset)
    at $R\!=\!4$ (left) and the same configuration extended to
    $R\!=\!6$ (right).  Each line is one of five seeds; per-round
    eval uses 20 episodes and is noisy.  The stars at the right
    edge mark the final \emph{deterministic 50-episode eval},
    coloured by the final verdict (green: $\bar C \le 5$, red:
    $\bar C > 5$).
    Three of five seeds in each run land cleanly inside the GO band
    while two settle at $\bar C\!\ge\!19$ — the same 3/5 split
    occurs in both panels, with different seed identities.
    Extending from $R\!=\!4$ to $R\!=\!6$ does not rescue the
    NO-GO seeds: more rounds do not change which basin a seed
    converges to.  This is the §D.4 ``two fixed points''
    diagnosis — the iteration converges, but to either of two
    attractors (safe versus over-conservative), determined by the
    initial PPO trajectory.%
  }
  \label{fig:vii_d_iicf_rounds}
\end{figure*}

%
% Cross-references: §VII-D §D.3 (results table) and §D.4
% (diagnosis — two-fixed-point interpretation).
% =====================================================================
\begin{figure*}[t]
  \centering
  \includegraphics[width=\textwidth]%
    {figures/fig_vii_d_iicf_rounds/fig_vii_d_iicf_rounds.pdf}
  \caption{%
    \textbf{IICF round trajectories — the bimodal final outcome
    persists across $R \in \{4, 6\}$.}
    Per-seed episode cost over the outer-loop rounds for the v3b
    configuration (window depth $W\!=\!2$, per-round critic reset)
    at $R\!=\!4$ (left) and the same configuration extended to
    $R\!=\!6$ (right).  Each line is one of five seeds; per-round
    eval uses 20 episodes and is noisy.  The stars at the right
    edge mark the final \emph{deterministic 50-episode eval},
    coloured by the final verdict (green: $\bar C \le 5$, red:
    $\bar C > 5$).
    Three of five seeds in each run land cleanly inside the GO band
    while two settle at $\bar C\!\ge\!19$ — the same 3/5 split
    occurs in both panels, with different seed identities.
    Extending from $R\!=\!4$ to $R\!=\!6$ does not rescue the
    NO-GO seeds: more rounds do not change which basin a seed
    converges to.  This is the §D.4 ``two fixed points''
    diagnosis — the iteration converges, but to either of two
    attractors (safe versus over-conservative), determined by the
    initial PPO trajectory.%
  }
  \label{fig:vii_d_iicf_rounds}
\end{figure*}
